% Step 7 historical note: exact orbit-volume model calculations
% S0910A03 Agent2 (Constructor) — Explore
%
% CORRECTED IN STEP 8. The exact identities below remain useful only under
% their stated symmetry and base hypotheses. The total-deficit transport,
% capacity, exhaustive-collapse, unconditional interval nil^3, and bounded
% cusp-length statements are superseded by
% step8_cover_mean_transport_and_model_scope_corrections.md.
% No fixed-topology pinching gap follows from this note.
%
% Sources: Lott (arxiv:1708.06780) eq. (2.9), Prop. 2.10;
%   Cheeger--Naber (Ann. Math. 2015); Naber--Zhang (Geom. Topol. 2016).

\section{The orbit-volume equation for negative Einstein metrics}

Let $M^{n+N}$ carry an adapted metric with free local isometric
$\mathbb R^N$-action:
\[
 \bar g=\sum_{I,J}G_{IJ}(dx^I+A^I)(dx^J+A^J)+g_B,
\]
where $G_{IJ}$ is a positive-definite symmetric matrix on the base $B^n$,
$A^I$ are connection one-forms, and $g_B$ is a Riemannian metric on $B$.
Put $v=\sqrt{\det G}$ (the orbit-volume function) and
$|F|^2=g^{\alpha\gamma}g^{\beta\delta}G_{IJ}F_{\alpha\beta}^I F_{\gamma\delta}^J$.

\textbf{Lott's equation} \cite[eq.\ (2.9)]{Lott2020}:
if $\operatorname{Ric}(\bar g)=\lambda\bar g$, then
\begin{equation}\label{eq:lott}
 \Delta_B v=\Bigl(\frac{|F|^2}{4}-\lambda N\Bigr)v.
\end{equation}
(The Laplacian is the geometer's $\Delta=\operatorname{tr}\nabla^2$, so
$\Delta v\ge0$ means $v$ is subharmonic.)

For the \textbf{negative normalized branch} $\lambda=-1$,
\begin{equation}\label{eq:negein}
 \Delta_B v=\Bigl(\frac{|F|^2}{4}+N\Bigr)v.
\end{equation}
Since $|F|^2\ge0$ and $N\ge1$, the coefficient is
$\ge N\ge1>0$.  Therefore:
\begin{quote}
\emph{Wherever $v>0$, the orbit-volume function $v$ is strictly
subharmonic.}
\end{quote}

This is qualitatively different from the Ricci-flat case $\lambda=0$,
where $\Delta v=(|F|^2/4)v\ge0$ merely forces $v$ constant and $F=0$
on parabolic bases, permitting nontrivial semiflat K\"ahler collapse
(K3 examples exist).

\section{Exclusion of free abelian negative Einstein collapse}

\textbf{Proposition.}  Let $(B^n,g_B)$ be a connected Riemannian manifold
(or orbifold), and let $\mathcal X$ be a Riemannian groupoid over $B$
with free abelian $\mathbb R^N$ symmetry ($N\ge1$) and
$\operatorname{Ric}_{\mathcal X}=\lambda g_{\mathcal X}$ with $\lambda<0$.
If $B$ is parabolic and $\det G$ is bounded above, then the free abelian
structure is impossible: no such $\mathcal X$ exists.

\textit{Proof.}
From \eqref{eq:negein}, $v=\sqrt{\det G}\ge0$ satisfies
$\Delta v\ge(-\lambda N)v$ with $-\lambda N>0$.
On the open set $\{v>0\}$, $v$ is strictly subharmonic.
Since $v\ge0$ everywhere and $\det G$ bounded above implies $v$ bounded
above, on a parabolic base a bounded-above subharmonic function is
constant (see \cite[Thm.\ 5.1]{Grigoryan1999}).
If $v=c>0$ is constant, then $\Delta v=0$ but the right side
$(-\lambda N)c>0$: contradiction.
If $v\equiv0$, then $\det G\equiv0$, contradicting the free abelian
structure (which requires positive-definite $G$). \qed

\subsection*{Scope in four dimensions}
For an exact adapted model, the proposition applies in any orbit dimension
provided its hypotheses are independently verified.  In base dimensions two
or three, a punctured compact smooth/orbifold base with polar punctures is a
standard parabolic case, but finite total four-volume does \emph{not} imply
that $v$ is bounded.  Thus no general free-circle or free-$T^2$ collapse
exclusion follows from volume alone.  For a circle base, periodic integration
excludes the exact abelian model.  For an interval, endpoint behavior must be
specified.  If the abelian action is transitive ($N=4$), the adapted metric is
left-invariant on an abelian group and hence flat, so it cannot satisfy
$\operatorname{Ric}=-g$.
\section{The quantitative near-Einstein identity under an exact adapted action}

Suppose the actual metric $g_i$ itself carries an exact local free isometric
$\mathbb R^N$-action and adapted decomposition.  This is an additional model
hypothesis; it is not supplied by the pinching assumption or by an unverified
Cheeger--Tian regular region.
If the region carries a local free $\mathbb R^N$-action adapted to $g_i$,
the vertical trace of $\operatorname{Ric}_{g_i}=-g_i+E_i$ gives
\[
 N\lambda_i+\operatorname{tr}_{\rm vert}E_i
 =-\frac12 g^{\alpha\beta}G^{IJ}G_{IJ;\alpha\beta}+\cdots+\frac{|F|^2}{4}
\]
with $\lambda_i=-1$ and $|\operatorname{tr}_{\rm vert}E_i|\le N\delta_i$.
Hence the near-Einstein orbit-volume inequality is
\begin{equation}\label{eq:nearein}
 \Delta_B v\ge\Bigl(\frac{|F|^2}{4}+N(1-\delta_i)\Bigr)v
 \ge N(1-\delta_i)\,v.
\end{equation}
For $\delta_i<1$, the coefficient $c_i=N(1-\delta_i)>0$.

\subsection*{Capacity estimate (RETRACTED)}
\emph{The capacity estimate ``$c\int\eta^2\le\int|\nabla\eta|^2$''
claimed in an earlier draft is false.  The IPP/AM--GM derivation has
a sign error: the squared term in the expansion has the wrong direction
to absorb the cross term.  Counterexample: $v=e^{\sqrt c\,x}$ on
$\mathbb R$ satisfies $v''=cv$ but $c\int\eta_R^2\asymp R$ while
$\int|\eta_R'|^2=O(1)$.  Only those parabolic or periodic arguments whose
separate global hypotheses are verified remain valid.}

\section{Weighted volume and corrected mean transport}

This historical calculation motivates the issue; the rigorous normal-cover
mean lemma and its hypotheses are in the Step~8 correction.

\subsection*{Volume pushdown}
On a regular abelian region, the total-space volume element pushes down
to the base as
\[
 d\mu_M=v\cdot(\text{lattice factors})\cdot d\mu_B.
\]
If $\mathcal F$ is a nonneg pointwise scalar on $M$ (e.g., a Ricci
or curvature defect), its integral over a tube $\pi^{-1}(U)$ is
\[
 \int_{\pi^{-1}(U)}\mathcal F\,d\mu_M
 =\int_U\Bigl(\int_{\text{fiber}}\mathcal F\Bigr)\,v\,d\mu_B
 =\int_U\overline{\mathcal F}\cdot v\,d\mu_B,
\]
where $\overline{\mathcal F}$ denotes the fiberwise average.

\subsection*{Deficit amplification under unwrapping}
When the torus fiber of length $\sim v^{1/N}$ is unwrapped by a local
normal cover $\hat\pi:\hat M\to M$, the covering degree is $\sim1/v$
(in appropriate units).  A noncollapsed ball $\hat B$ on the cover
satisfies
\[
 \operatorname{Vol}(\hat B)\sim1
 \quad(\text{noncollapsed}).
\]
The pullback of $\mathcal F$ satisfies
\[
 \fint_{\hat B}\hat\pi^*\mathcal F
 \sim\frac{1}{v}\cdot\fint_{B}\mathcal F\cdot v\,d\mu_B.
\]
Thus \textbf{small $v$ amplifies pointwise defects on the cover}, as the
Strategist noted: a region with tiny fiber volume may carry little
base-space error ($\overline{\mathcal F}\cdot v$ small) yet substantial
error on the noncollapsed cover chart ($\overline{\mathcal F}$ itself
large).

\subsection*{Correct object: a local mean}
The covering-degree identity does \emph{not} imply that the cover integral is
bounded by the total downstairs deficit: the degree may diverge.  For complete
finite covers, numerator and volume both multiply, preserving the mean; for
cover balls, multiplicity is nonconstant.  Step~8 proves a normal-cover
mean-transfer lemma using disjoint deck translates and Bishop--Gromov
comparison.  Absolute global integral smallness alone is insufficient.

\subsection*{No capacity consequence}
The retracted capacity inequality gives no volume or boundary-cost estimate.
Subharmonicity retains only the conditional Liouville statement for an exact
adapted action over a parabolic base with bounded $v$.  It does not classify
the collapsed regime or force the surviving geometry to be cusp-like.

\section{Why these model identities do not yet apply to the pinched sequence}

Fixing a critical Ishida--LeBrun manifold gives absolute calibrated deficit
convergence and, by Step~6, everywhere unit-scale collapse for a hypothetical
sequence.  It does \emph{not} give the Einstein-only Cheeger--Tian curvature
decomposition asserted in an earlier draft: an $L^\infty$ bound on
$\operatorname{Ric}+g$ does not control its derivatives and therefore does not
produce bounded-curvature CFG regions or a smooth Einstein groupoid by that
route.

Even if an exact adapted abelian limit model were available, the parabolic
argument would still require a genuinely parabolic base with polar ends and a
bounded orbit-volume function; finite total volume supplies only an $L^1$
condition.  Hence the exact abelian model theorem cannot be counted as an
exclusion of the actual collapsed sequence.

The valid replacement is the Step~8 cover-mean theorem.  At points satisfying
uniform Naber--Zhang low-dimensional GH control, maximal fibered nilpotency
rank, and a nondegenerating model scale, the local universal cover has fixed
harmonic radius.  A sheet-independent mean-transfer inequality then carries a
small base-ball calibrated mean to the cover chart, where the forbidden local
Einstein spectrum gives a contradiction.  What remains open is the volume
occupied by points failing these uniform Naber--Zhang hypotheses.

\section{The nil$^3$ case: corrected analysis}

For 4D collapse with $\operatorname{nil}^3$ (Heisenberg) local symmetry,
the base is 1-dimensional and the fiber is a nilmanifold with
$[e_1,e_2]=e_3$.  From the hypersurface Ricci formula
and the Milnor fiber eigenvalues $r_1=r_2=-\mu$, $r_3=+\mu$
where $\mu=a_3^2/(2a_1^2a_2^2)$, the orbit-volume equation becomes
\begin{equation}\label{eq:nil3}
 v''=\Bigl(3-\mu\Bigr)v
 =\Bigl(3-\frac{a_3^2}{2a_1^2a_2^2}\Bigr)v.
\end{equation}
The Heisenberg structure constant \textbf{reduces} the coefficient from~3,
because the fiber scalar curvature $-\mu$ is negative.  The competing
$3+\mu$ formula in an earlier peer message had swapped the horizontal and
central fiber-Ricci signs; the peer has since retracted it.  The
$\mathbb{CH}^2$ cusp and \texttt{verify\_nil3\_sign\_check.py} independently
verify the displayed sign.

\subsection*{Circle base exclusion via maximum principle}
The full Einstein ODE system includes
$\rho'=2\rho D$ where $D=f_3-f_1-f_2$ and $\rho=2\mu$.
At an interior maximum of $\rho$, $D=0$ and
$\rho''=2\rho D'=\rho(3\rho-2)$.  Since $\rho''\le0$ at a maximum,
$\rho\le2/3$.  On a circle, $\rho$ attains its maximum in the interior,
so $\rho\le2/3$ everywhere, hence $\mu\le1/3$ and
\[
 v''\ge\frac83v>0.
\]
Integrating: $0=\int v''\ge\frac83\int v>0$, contradiction.

\subsection*{Interval base: cusps exist}
On a half-line, the standard $\mathbb{CH}^2$ horospherical metric
\[
 g=ds^2+e^{2s/\sqrt6}(\theta_1^2+\theta_2^2)
   +\tfrac23e^{4s/\sqrt6}\theta_3^2,
 \qquad\mu=\tfrac13,
\]
satisfies $\operatorname{Ric}=-g$ with $v\sim e^{4s/\sqrt6}$.
Taking a lattice quotient in the Heisenberg fiber gives a finite-volume
nil cusp.  This is a genuine negative Einstein solution with
$\text{nil}^3$ structure.  A two-ended finite interval is excluded by
convexity only when the endpoint behavior controls $\rho=2\mu$ (for example
$\limsup_{\partial I}\rho<6$) and $v\to0$ at both finite-distance endpoints.
If $\rho$ diverges at an endpoint, the interior maximum argument gives no
global convexity estimate.  A cusp can have infinite base length and finite
volume; see the Step~8 correction.

\subsection*{Unresolved singular and inaccessible regimes}
The model equations do not control rank-drop strata, singular fibers,
residual four-dimensional cores, curvature concentration, or regions whose
local covers become regular only at scales tending to zero.  Chern--Gauss--Bonnet
gives a fixed total $L^2$ curvature budget on fixed topology, but bounded Ricci
alone does not decompose the manifold into finitely many blowup balls and a
bounded-curvature complement.  The supported Step~8 statement instead excludes
positive-volume families of points satisfying uniform Naber--Zhang cover
hypotheses.  Determining the volume occupied by the complementary failure modes
is open.

\subsection*{Updated summary of the four-dimensional negative critical route}
\begin{center}
\begin{tabular}{ll}
\textbf{Regime} & \textbf{Status} \\
\hline
Noncollapsed regular patches & Excluded (Step~6 local spectral contradiction) \\
Exact abelian model, parabolic base & Excluded if $v$ bounded and action adapted \\
Exact abelian model, circle base & Excluded by periodic integration \\
Exact abelian model, interval base & Requires endpoint data \\
Exact $\text{nil}^3$ model, circle base & Excluded ($\rho\le2/3$, $v''\ge\frac83v$) \\
Exact $\text{nil}^3$ model, finite interval & Excluded only with endpoint $\limsup\rho<6$ and $v\to0$ \\
$\mathbb{CH}^2$ nil cusp & Exists; finite volume and infinite base length \\
Actual collapsed near-Einstein sequence & Open; exact models do not exhaust it \\
Positive branch, all dimensions & Open (separate mechanisms) \\
\end{tabular}
\end{center}

These are exact-model statements relevant to the four-dimensional negative
branch.  They do not classify a fixed critical Ishida--LeBrun near-Einstein
sequence and do not establish a pinching gap.

\medskip
\noindent\textbf{Key files:}
\texttt{step6\_local\_projector\_obstruction\_and\_collapse\_reduction.tex}
(local spectral exclusion),
\texttt{step6\_collapsed\_unwrapping\_route.md}
(collapse route plan),
\texttt{scripts/verify\_orbit\_volume\_negative\_einstein.py}
(algebraic verification).
This note:
\texttt{step7\_orbit\_volume\_negative\_einstein\_collapse.tex}.
\end{LaTeX}
